% !TEX root = ../libretto.tex
% Sezione: Quaresima e Pasqua
\momento{Quaresima e Pasqua}
\begin{songs}{titoli}

% Fonte: LibrettoAccordi.pdf (Canti Coro), O14; testo e struttura confermati dall'utente.
% Nell'elenco: Ceneri (l'elenco dice Buttazzo, ma l'utente conferma RnS).
\beginsong{Abbà misericordia}[by={RnS}]
\beginverse
Non sono degno di essere qui, ho abbandonato la tua casa,
ho dissipato i tuoi beni,
Padre, ho peccato contro il cielo e contro di Te.
\endverse
\beginchorus
Abbà, misericordia, Abbà, misericordia, Abbà, Abbà.
Abbà, misericordia, Abbà, misericordia, Abbà, Abbà, Abbà.
\endchorus
\beginverse
Non sono degno di esser tuo figlio, in Gesù Cristo rialzami
e in Lui ridammi la vita,
Padre, ho peccato contro il cielo e contro di Te.
\endverse
\beginverse
Non sono degno del tuo amore, riempi il mio cuore del tuo Spirito,
insieme a Te farò festa per sempre,
Padre, ho peccato contro il cielo e contro di Te.
\endverse
\endsong

% Fonte: LibrettoAccordi.pdf (Canti Coro), E8.
% Nell'elenco: Alleluia, Pasqua.
\beginsong{Alleluia, Cristo è risorto veramente}[by={F. Buttazzo}]
\beginchorus
Alleluia! Alleluja! Alleluja! Alleluja!
Cristo è risorto veramente. Alleluja! Alleluja! \rep{2}
\endchorus
\beginverse
Cristo ha inaugurato il suo regno,
salvezza e vita sono con lui.
Rallegriamoci, esultiamo con lui:
la morte è stata vinta, la vita ora trionfa.
\endverse
\endsong

% Fonte: "99 - Raccolte e libretti accordi/Il mistero pasquale_ gv.pdf" (pag. 77).
% Nell'elenco: Comunione, Settimana Santa. CD: Il mistero pasquale (2011).
\beginsong{Amore abbandonato}[by={Gen Verde; testo P. A. Sequeri (CD Il mistero pasquale)}]
\beginverse
Giorno verrà
del tuo ritorno in festa:
un vino nuovo porterai
ed io berrò con te.
Come il tralcio resto in te
e accolgo in me la vita.
Dal tuo ferito amore io
non mi separerò.
\endverse
\beginverse
Canto per te
che doni la tua vita:
nel sangue mi purifichi,
agnello immacolato.
Del mio cuore arido
tu spezzi le catene.
Più nulla ci separerà:
un cuore nuovo avrò.
\endverse
\beginverse
Resta con me
la sera della vita:
nel segno della croce io
ti riconoscerò.
Nel mio amore fragile
tu accendi la speranza.
Amore abbandonato, tu
non mi abbandonerai.
\endverse
\endsong

% Fonte: LibrettoAccordi.pdf (Canti Coro), O18.
% Nell'elenco: Finale, Settimana Santa.
\beginsong{Ave verum}[by={Gen Verde (CD Il mistero pasquale)}]
\beginchorus
Ave, verum corpus,
natum de Maria virgine,
vere passum immolatum
in cruce pro homine.
Cujus latus perforatum
unda fluxit et sanguine,
esto nobis praegustatum
in mortis examine.
O Iesu dulcis, o Iesu pie,
o Iesu, fili Mariae,
miserere, miserere, miserere mei.
\endchorus
\beginverse*
Amen. Amen. Amen. Amen.
\endverse
\endsong

% Fonte: LibrettoAccordi.pdf (Canti Coro), O21.
% Nell'elenco: Pasqua.
\beginsong{Canto dell'acqua}[by={Gen Verde (CD Il mistero pasquale)}]
\beginverse
L'acqua che sommerge il buio delle colpe scorre qui.
L'acqua che ci rende il cuore trasparente scorre qui.
È zampillo nato da roccia, è sorgente di fecondità,
è fontana del tuo giardino, onda di libertà.
L'acqua che ci immerge dentro la tua morte scorre qui.
L'acqua che ci dà la vita che risorge scorre qui.
È ruscello nato da piaga, fresco fuoco che ci inonderà,
acqua viva e bianca rugiada, goccia d'eternità.
L'acqua viva, grembo della vita nuova, scorre qui.
\endverse
\endsong

% Fonte: "15 - Quaresima, Triduo, Pasqua/come è grande il tuo amore.pdf" (foglio accordi).
% Nell'elenco: Ingresso, Quaresima (autore in elenco: Ecco perché canto).
\beginsong{Com'è grande il tuo amore}[by={A. Benedetti}]
\beginverse
Siamo qui
per cercare il tuo volto, o Dio,
tu, sorgente di vita sei,
acqua viva che sgorga dalla roccia.
Solo in te
ritroviamo il perdono e ci rialziamo di nuovo
se ci tendi la tua mano ancora, Gesù.
\endverse
\beginchorus
Com'è grande il tuo amore
per tutte le creature,
tu ci ami, Signore,
e non ci abbandoni mai. \rep{2}
\endchorus
\beginverse
Su di noi
hai posato lo sguardo, o Dio,
ci hai chiamati per nome, tu
sei la luce che illumina la notte.
Solo in te
ritroviamo il perdono e ci rialziamo di nuovo
se ci tendi la tua mano ancora, Gesù.
\endverse
\beginverse*
E non ci abbandoni mai,
tu non ci abbandoni.
\endverse
\endsong

% Fonte: LibrettoAccordi.pdf (Canti Coro), O2.
% Nell'elenco: Lavanda, Settimana Santa.
\beginsong{Come io ho amato voi}[by={A. Porfiri}]
\beginverse
Amatevi a vicenda come io ho amato voi,
così dice il Signore:
niente divisioni, liti o discordie,
l'amore regni tra voi.
\endverse
\beginverse
Servitevi a vicenda come io ho servito voi,
così dice il Signore:
chi di voi è primo ultimo sarà
e l'amore regni tra voi.
\endverse
\beginverse
Donatevi a vicenda come io ho fatto a voi,
così dice il Signore:
siete un solo corpo, Popolo di Dio,
l'amore regni tra voi.
\endverse
\endsong

% Fonte: LibrettoAccordi.pdf (Canti Coro), G17.
% Nell'elenco: Offertorio, Pasqua.
\beginsong{Come un fiume}[by={Gen Verde (CD Il mistero pasquale)}]
\beginverse
Come un fiume il frutto della tua vittoria
viene a Te.
Tutto quello che ci hai dato nella storia
diamo a Te.
Ecco il mondo, immensa patena
dei fermenti dell'umanità;
ecco il pane ardente dei giorni
che il tuo amore ci dà.
Ecco il mondo, calice immenso
d'ogni pena, d'ogni povertà;
Ecco il vino ardente del tempo,
questo tempo che va.
\endverse
\endsong

% Fonte: LibrettoAccordi.pdf (Canti Coro), O13.
% Nell'elenco: Ingresso, Ceneri.
\beginsong{Dal profondo, o Signore}[by={F. Buttazzo}]
\beginchorus
Dal profondo, o Signore, gridiamo il tuo nome:
dalle strade del mondo vogliamo tornare a te.
Dal profondo, o Signore, gridiamo il tuo nome:
solo in te la vita risplenderà.
\endchorus
\beginverse
E tu ci guiderai sulla via che porta a te;
donerai salvezza all'umanità.
\endverse
\beginverse
E tu accoglierai l'uomo che ritorna a te;
una grande festa tu farai per lui.
\endverse
\beginverse
E tu illuminerai ogni nostra oscurità;
splenderà il tuo volto dentro di noi.
\endverse
\beginverse
E tu ridonerai la tua pace all'umanità;
e verrà il tuo Regno in mezzo a noi.
\endverse
\endsong

% Fonte: LibrettoAccordi.pdf (Canti Coro), O12.
% Nell'elenco: Offertorio, Settimana Santa.
\beginsong{Dayenu Adonai}[by={Gen Verde (CD Il mistero pasquale)}]
\beginverse
Se dopo averci fatti uscire dall'Egitto
Tu non ci avessi sostenuti con la manna,
dayenu, dayenu Adonai,
dayenu, dayenu Adonai.
Se dopo averci sostenuti con la manna
Tu non ci avessi consegnato la tua Legge,
dayenu, dayenu Adonai,
dayenu, dayenu Adonai.
Se dopo averci consegnato la tua Legge
Tu non ci avessi fatti entrare in Israele,
dayenu, dayenu Adonai,
ci sarebbe bastato, Signore.
\endverse
\beginverse
Se dopo averci definito tuoi amici
Tu non ti fossi inginocchiato a lavarci,
dayenu, dayenu Adonai,
dayenu, dayenu Adonai.
Se quando Tu ti sei chinato a lavarci
poi non ti fossi a noi spezzato come pane,
dayenu, dayenu Adonai,
dayenu, dayenu Adonai.
Se quando Tu ti sei spezzato come pane
poi non ti fossi consegnato all'abbandono
dayenu, dayenu Adonai,
ci sarebbe bastato, Signore.
\endverse
\beginverse
Per tutto questo cosa mai potrò donarti?
Se ti rendessi oro e perle senza pari,
non potrà mai bastare, Adonai,
o Signore, non basterà mai.
Se ti rendessi oro e perle senza pari
ed aggiungessi il coro immenso dei respiri,
non potrà mai bastare, Adonai,
o Signore, non basterà mai.
Se ti rendessi il coro immenso dei respiri
ed aggiungessi terre, cieli ed universi,
non potrà mai bastare, Adonai,
o Signore non basterà mai.
\endverse
\beginverse
Noi ti portiamo tralci spogli e solchi vuoti,
Tu ci darai stille d'uva e farina;
solo Tu, solo Tu Adonai,
solo Tu, solo Tu basterai.
Noi ti portiamo stille d'uva e farina,
Tu ne farai il tuo corpo, o Signore;
solo Tu, solo Tu Adonai,
solo Tu, solo Tu basterai.
Noi ti portiamo mille vite e mille cuori,
Tu ci farai un solo corpo, un solo cuore;
solo Tu, solo Tu Adonai,
solo Tu, solo Tu basterai.
Solo Tu, solo Tu basterai.
\endverse
\endsong

% Fonte: LibrettoAccordi.pdf (Canti Coro), J2.
% Nell'elenco: Comunione, Pasqua.
\beginsong{E sei rimasto qui}[by={Gen Rosso}]
\beginverse
Perché la sete d'infinito? Perché la fame d'immortalità?
Sei tu che hai messo dentro il cuore il desiderio dell'eternità.
Ma Tu sapevi che quel vuoto lo colmavi tu,
per questo sei venuto in mezzo a noi.
\endverse
\beginchorus
E sei rimasto qui visibile mistero.
E sei rimasto qui cuore del mondo intero.
E rimarrai con noi finché questo universo girerà, salvezza dell'umanità!
\endchorus
\beginverse
Si apre il cielo del futuro, il muro della morte ormai non c'è
Tu, Pane Vivo, ci fai Uno, richiami tutti i figli attorno a te.
E doni il Tuo Spirito che lascia dentro noi
il germe della sua immortalità.
\endverse
\beginverse
Presenza vera nel mistero, ma più reale di ogni realtà.
Da te ogni cosa prende vita, e tutto un giorno a te ritornerà.
Varcando l'infinito tutti troveremo in Te
un Sole immenso di felicità.
\endverse
\endsong

% Fonte: LibrettoAccordi.pdf (Canti Coro), O17.
% Nell'elenco: Ingresso, Settimana Santa.
\beginsong{È venuta l'ora}[by={Gen Verde (CD Il mistero pasquale)}]
\beginverse
È venuta l'ora; c'è pianto e festa stasera,
ci raccogli insieme con te e ti consegni a noi, Signore.
\endverse
\beginchorus
Kyrie eleison, Christe eleison, Kyrie eleison, eleison.
\endchorus
\beginverse
È venuta l'ora; tu, come servo, Signore
versi un'acqua di libertà su piedi rossi di catene.
\endverse
\beginverse
È venuta l'ora d'amore senza misura;
tutto il tuo mistero ci dai, che sa di cielo e sa di grano.
\endverse
\beginverse
La tua gloria ha il colore delle piaghe,
la tua croce di vittoria ci rivestirà.
\endverse
\endsong

% Fonte: LibrettoAccordi.pdf (Canti Coro), K1.
% Nell'elenco: Finale, Pasqua.
\beginsong{Esulta il cielo}[by={F. Buttazzo}]
\beginchorus
Esulta il cielo e canta a noi la novità:
Gesù è risorto dalla morte!
Esulta l'uomo per la vita che gli donerà
Gesù, il Signore.
Esulta il cielo e canta a noi la novità:
Gesù è risorto dalla morte!
E si spalancano le porte dell'eternità
per tutti noi.
\endchorus
\beginverse
Il vero Agnello che toglie i peccati è lui,
lui che morendo ha distrutto la morte,
e risorgendo ha ridato la vita,
vita per l'eternità.
\endverse
\beginverse
In lui risplende la vita del Padre,
a noi è donata la vita dei figli.
E sono aperte le porte del Regno
per chi crede nel suo nome.
\endverse
\beginverse
In lui, che ha vinto il peccato e la morte,
l'universo si rinnova e risorge;
l'uomo ritorna alle vere sorgenti
della vita insieme a Dio.
\endverse
\beginverse
Ora continua ad offrirsi per tutti noi,
lui ci difende, è nostro «avvocato».
L'uomo che vive amando i fratelli
la salvezza abbraccerà.
\endverse
\endsong

% Fonte: LibrettoAccordi.pdf (Canti Coro), O5.
% Nell'elenco: Settimana Santa.
\beginsong{Il Signore è la luce}[by={M. Giombini}]
\beginverse
Il Signore è la luce che vince la notte.
Gloria, gloria! Cantiamo al Signore!
\endverse
\beginverse
Il Signore è il coraggio che vince il terrore.
Gloria, gloria! Cantiamo al Signore!
\endverse
\beginverse
Il Signore è il sereno che vince la pioggia.
Gloria, gloria! Cantiamo al Signore!
\endverse
\beginverse
Il Signore è la vita che vince la morte.
Gloria, gloria! Cantiamo al Signore!
\endverse
\beginverse
Il Signore è l'amore che vince il peccato.
Gloria, gloria! Cantiamo al Signore!
\endverse
\beginverse
Il Signore è la gioia che vince l'angoscia.
Gloria, gloria! Cantiamo al Signore!
\endverse
\beginverse
Il Signore è la pace che vince la guerra.
Gloria, gloria! Cantiamo al Signore!
\endverse
\endsong

% Fonte: LibrettoAccordi.pdf (Canti Coro), O10.
% Nell'elenco: Lavanda, Settimana Santa.
\beginsong{Io vi do un grande esempio}[by={D. Haas}]
\beginchorus
Io vi do un grande esempio, con amore lavo i vostri piedi,
se ora tutti voi farete come me l'amore crescerà.
\endchorus
\beginverse
Ora, Pietro, il mio gesto non capisci,
un giorno non lontano capirai
cos'è l'amore.
\endverse
\beginverse
Ti stupisce quello che io sto facendo,
io mi chino per lavare i tuoi piedi: questo è l'amore.
\endverse
\beginverse
Guarda e impara da me.
Chi lava i piedi all'altro capirà
cos'è l'amore.
\endverse
\beginverse
Questo è l'amore più grande:
dare con gioia la tua vita
per un amico.
\endverse
\beginverse
Vai e vivi nell'amore;
tu sarai mio discepolo fedele
se amerai.
\endverse
\beginverse
Sono io che ho scelto te.
Io ti mando ad annunciare il mio regno
ai miei fratelli.
\endverse
\endsong

% Fonte: "99 - Raccolte e libretti accordi/Il mistero pasquale_ gv.pdf" (pag. 69).
% Nell'elenco: Adorazione della Croce, Settimana Santa. CD: Il mistero pasquale (2011). In Quaresima e Pasqua su indicazione dell'utente (non tra i Mariani).
\beginsong{Lì con te}[by={Gen Verde; testo P. Stradi (CD Il mistero pasquale)}]
\beginverse
Figlio mio, tu sei qui;
sanguina ancora il tuo perché.
Scende la spada nel cuore,
scende sul mondo la notte;
qui sul mio grembo,
come agnello svenato sei qui.
\endverse
\beginverse
Sono qui, qui con te.
Niente può farti male, ormai;
non urla più l'uragano,
tutto lo strazio è lontano;
qui sul mio grembo,
per un'ultima volta sei qui.
\endverse
\beginchorus
Ed ora ti porteranno via
strappandoti da me,
ma qui sul mio grembo resterà
la tua impronta, Figlio,
ed il sangue tuo su di me.
\endchorus
\beginverse
Ora addio, Figlio mio.
L'ultimo bacio che ti do
ti porta tutto l'amore
di chi non trova parole
per dirti grazie
e per chiedere la tua pietà.
\endverse
\beginchorus
Ed ora ti porteranno via
strappandoti da me,
ma Figlio, non t'abbandonerò
e davanti al sepolcro sbarrato, lì io starò.
\endchorus
\beginchorus
Ed ora ti porteranno via;
in piedi aspetterò
finché da quel grembo nascerai
e per sempre vivo mi porterai
lì con te.
\endchorus
\endsong

% Fonte: LibrettoAccordi.pdf (Canti Coro), O11.
% Nell'elenco: Ingresso, Settimana Santa.
\beginsong{L'ultima cena con i suoi}[by={F. Buttazzo}]
\beginverse
È vicina la Pasqua, festa degli Ebrei,
e Gesù siede a cena coi discepoli suoi.
Raccontava la storia dell'Egitto e di Mosè,
e, spezzando il pane, disse loro così.
\endverse
\beginchorus
Prendete e mangiatene, amici miei: il pane è il mio corpo, donato a voi.
Bevete il mio calice e fate questo in memoria di me.
\endchorus
\beginverse
Quella sera si apriva il cuore di Gesù,
fece un dono speciale ai discepoli suoi:
anche loro saranno sacerdoti come Gesù,
sull'altare diranno ciò che ha detto Gesù.
\endverse
\beginverse
In quell'ultima cena con i Dodici
Lui veniva tradito da un discepolo.
Con immenso dolore anche a Giuda il pane donò
e con vero amore disse loro così.
\endverse
\endsong

% Fonte: "99 - Raccolte e libretti accordi/Il mistero pasquale_ gv.pdf" (pag. 20).
% Nell'elenco: Lavanda, Settimana Santa. CD: Il mistero pasquale (2011).
\beginsong{Non c'è amore più grande}[by={Gen Verde (CD Il mistero pasquale)}]
\beginchorus
Non c'è amore più grande di chi dà la vita per i suoi.
Non c'è amore più grande: io do la mia vita per voi.
\endchorus
\beginverse
Il Signore si alzò da tavola,
versò dell'acqua in un catino,
e cominciò a lavare i piedi ai discepoli:
ad essi volle lasciare questo esempio.
\endverse
\beginverse
«Quello che io faccio, ora non lo comprendi,
ma lo comprenderai un giorno».
«Signore, tu lavi i piedi a me?».
Gesù gli rispose dicendo:
«Se non ti laverò, non avrai parte con me».
\endverse
\beginverse
«Se vi ho lavato i piedi,
io, Signore e Maestro,
quanto più voi avete il dovere
di lavarvi i piedi l'un l'altro».
\endverse
\beginverse
Gesù disse ai suoi discepoli:
«Da questo tutti sapranno
che siete miei discepoli,
se vi amerete gli uni gli altri».
\endverse
\beginverse
«Vi do un comandamento nuovo:
che vi amiate gli uni gli altri
come io ho amato voi»,
dice il Signore.
\endverse
\beginverse
Fede, speranza e carità,
tutte e tre rimangano tra voi:
ma più grande di tutte è la carità.
\endverse
\beginchorus
Non c'è amore più grande di chi dà la vita per i suoi.
Non c'è amore più grande: io do la mia vita per voi, per voi.
\endchorus
\endsong

% Fonte: LibrettoAccordi.pdf (Canti Coro), O1.
% Nell'elenco: Ingresso, Settimana Santa.
\beginsong{Osanna al Figlio di David}[by={Canto popolare}]
\beginchorus
Osanna al figlio di David! Osanna al Redentor. \rep{2}
\endchorus
\beginverse
Apritevi o porte eterne, avanzi il Re della gloria;
adorin cielo e terra l'eterno suo poter.
\endverse
\beginverse
O monti stillate dolcezza, il Re d'amor s'avvicina;
si dona pane vivo ed offre pace al cuor.
\endverse
\beginverse
O Vergine presso l'Altissimo trovasti grazia e onor;
soccorri i tuoi figliuoli donando il Salvator.
\endverse
\beginverse
Verrai un giorno Giudice, o mite e buon Gesù;
rimetti i nostri falli nel tempo del perdon.
\endverse
\endsong

% Fonte: testo fornito dall'utente. Non presente nell'elenco. Autore indicato dall'utente.
% Domenica delle Palme. Il ritornello all'inizio e alla fine si canta due volte (qui scritto una volta).
\beginsong{Rallegrati, Gerusalemme}[by={P. Dargenio}]
\beginchorus
Rallegrati, Gerusalemme,
esulta tra le tue mura,
perché il Signore, adesso, viene da te
tra petali e mantelli, tra canti di gioia:
Osanna, osanna al Re!
\endchorus
\beginverse
Le folle degli ebrei cantavano a gran voce:
Osanna, osanna al Re!
Ti festeggiavano con palme nelle mani:
Osanna, osanna al Re!
\endverse
\beginverse
Le folle degli ebrei non possono tacere:
Osanna, osanna al Re!
Ma se tacessero le pietre griderebbero:
Osanna, osanna al Re!
\endverse
\beginverse
Le folle degli ebrei stendevano i mantelli:
Osanna, osanna al Re!
E ti acclamavano Messia d'Israele:
Osanna, osanna al Re!
\endverse
\beginverse*
Osanna, osanna al Re!
Osanna, osanna al Re!
\endverse
\endsong

% Fonte: LibrettoAccordi.pdf (Canti Coro), K3 (li' tra i canti finali).
% Nell'elenco: Finale, Pasqua. Omesso il finale "a cori scontrati".
\beginsong{Resurrezione}[by={Gen Rosso}]
\beginverse
Che gioia ci hai dato, Signore del cielo, Signore del grande universo!
Che gioia ci hai dato, vestito di luce, vestito di gloria infinita,
vestito di gloria infinita!
\endverse
\beginverse
Vederti risorto, vederti Signore, il cuore sta per impazzire!
Tu sei ritornato, tu sei qui tra noi, e adesso ti avremo per sempre,
e adesso ti avremo per sempre.
\endverse
\beginverse
Chi cercate, donne, quaggiù? Chi cercate, donne, quaggiù?
Quello che era morto non è qui:
è risorto, sì! Come aveva detto anche a voi, voi gridate a tutti che
è risorto Lui, a tutti che è risorto Lui!
\endverse
\beginverse
Tu hai vinto il mondo, Gesù! Tu hai vinto il mondo, Gesù!
Liberiamo la felicità!
E la morte, no, non esiste più, l'hai vinta tu e hai salvato tutti noi,
uomini con te, tutti noi, uomini con te.
\endverse
\endsong

% Fonte: LibrettoAccordi.pdf (Canti Coro), O8. C'e' anche "6 - Offertorio/SERVIRE è REGNARE.pdf".
% Nell'elenco: Lavanda, Settimana Santa. CD: Il mistero pasquale (2011).
\beginsong{Servire è regnare}[by={Gen Verde (CD Il mistero pasquale)}]
\beginverse
Guardiamo a te che sei Maestro e Signore:
chinato a terra stai, ci mostri che l'amore
è cingersi il grembiule, sapersi inginocchiare,
ci insegni che amare
è servire.
\endverse
\beginchorus
Fa' che impariamo, Signore, da te,
che il più grande è chi più sa servire,
chi si abbassa e chi si sa piegare,
perché grande è soltanto l'amore.
\endchorus
\beginverse
E ti vediamo poi, Maestro e Signore,
che lavi i piedi a noi che siamo tue creature;
e, cinto del grembiule, che è il manto tuo regale,
ci insegni che servire
è regnare.
\endverse
\endsong

% Fonte: LibrettoAccordi.pdf (Canti Coro), O6 (sezione Settimana Santa - Ceneri).
% Nell'elenco: senza tipologia/periodo.
\beginsong{Ti saluto, o croce santa}[by={Gazzera, Damilano}]
\beginchorus
Ti saluto, o croce santa, che portasti il Redentor:
gloria, lode, onor ti canta ogni lingua ed ogni cuor.
\endchorus
\beginverse
Sei vessillo glorioso di Cristo, sua vittoria e segno d'amor:
il suo sangue innocente fu visto come fiamma sgorgare dal cuor.
\endverse
\beginverse
Tu nascesti fra braccia amorose d'una Vergine madre, o Gesù;
tu moristi fra braccia pietose d'una croce che data ti fu.
\endverse
\beginverse
O Agnello divino, immolato sulla croce crudele, pietà!
Tu, che togli dal mondo il peccato, salva l'uomo che pace non ha.
\endverse
\beginverse
Dona a tutti speranza, Signore, crocifisso e risorto per noi:
tu che effondi la pace del cuore nel tuo Spirito di santità.
\endverse
\beginverse
Del giudizio nel giorno tremendo sulle nubi del cielo verrai.
Piangeranno le genti, vedendo qual trofeo di gloria sarai.
\endverse
\endsong

% Fonte: testo fornito dall'utente.
% Nell'elenco: Finale, Pasqua. CD: Il mistero pasquale (2011).
\beginsong{Tutto il mondo deve sapere}[by={Gen Verde (CD Il mistero pasquale)}]
\beginchorus
Alleluia, alleluia, alleluia. \rep{2}
\endchorus
\beginverse
Come potrò raccontare?
È una gioia che fa piangere e fa gridare.
Io l'ho visto con i miei occhi,
era vivo, era Lui.
E m'ha chiamata per nome,
era la sua voce, era il mio Signore.
Io l'ho visto con i miei occhi,
era lì davanti a me.
\endverse
\beginverse
Gli angeli ci hanno parlato
davanti a quel sepolcro spalancato:
«Donne, il Signore è risorto,
non cercatelo qui».
Che corsa senza respiro
per gridare a tutti di quel mattino.
Tutto il mondo deve sapere
che è rimasto qui tra noi.
\endverse
\beginverse
Alba di un tempo diverso,
è il mattino dei mattini per l'universo.
Tutto già profuma d'eterno,
c'è il Risorto fra noi.
E le sue piaghe e la gloria
sono vive dentro la nostra storia,
segni di un amore che resta
qui per sempre, qui con noi.
\endverse
\beginverse*
Alleluia.
\endverse
\endsong

% Fonte: LibrettoAccordi.pdf (Canti Coro), O9 (sezione Settimana Santa - Ceneri).
% Nell'elenco: senza tipologia/periodo.
\beginsong{Venite dal profondo}[by={R. J. Stamps}]
\beginverse
Venite dal profondo dei tempi, o cuori dalla sete piagati.
Aprite il cuore alla gioia più profonda.
\endverse
\beginchorus
Dio ci ha messo il suo Corpo tra le mani. \rep{3}
\endchorus
\beginverse
Stasera l'acqua si trasforma in vino, sul monte il pane è moltiplicato.
Precoce è il frutto della vigna in fiore.
\endverse
\beginverse
Come ferita aperta è la parola. Parola di un eterno senza fine,
il Verbo si fa carne e nutrimento.
\endverse
\endsong

\end{songs}
